\begin{animation}[id="houses-schools", label="Houses and Schools -- Interval DP with D&C Optimization"]
\shape{a}{Array}{size=6, data=[2, 3, 5, 1, 4, 6], labels="1..6", label="Children per house"}
\shape{dp}{DPTable}{rows=3, cols=7, row_labels="j=0..2", col_labels="i=0..6", label="dp[j][i]"}
\shape{cost_val}{Array}{size=1, data=[0], label="cost(l,r)"}

\step
\narrate{Houses and Schools: place k=2 schools among n=6 houses to minimize total walking distance. Children counts: [2, 3, 5, 1, 4, 6]. We build dp[j][i] = min cost using j schools for houses 1..i.}

\step
\apply{dp.cell[0][0]}{value=0}
\recolor{dp.cell[0][0]}{state=done}
\narrate{Base case: dp[0][0] = 0. Zero schools covering zero houses costs nothing. All other dp[0][i] are infinite (cannot cover houses with no schools).}

\step
\recolor{dp.cell[1][1]}{state=current}
\highlight{a.cell[0]}
\apply{dp.cell[1][1]}{value=0}
\apply{cost_val.cell[0]}{value=0}
\narrate{dp[1][1] = cost(1,1) = 0. One school at house 1 serves only house 1. No walking needed.}

\step
\recolor{dp.cell[1][1]}{state=done}
\recolor{dp.cell[1][2]}{state=current}
\highlight{a.range[0:1]}
\apply{dp.cell[1][2]}{value=2}
\apply{cost_val.cell[0]}{value=2}
\narrate{dp[1][2] = cost(1,2) = 2. One school for houses 1-2: weighted median at house 2 (weight 3 $\geq$ half of 5). Cost = 2*1 + 3*0 = 2.}

\step
\recolor{dp.cell[1][2]}{state=done}
\recolor{dp.cell[1][3]}{state=current}
\highlight{a.range[0:2]}
\apply{dp.cell[1][3]}{value=7}
\apply{cost_val.cell[0]}{value=7}
\narrate{dp[1][3] = cost(1,3) = 7. Weighted median at house 3. Cost = 2*2 + 3*1 + 5*0 = 7.}

\step
\recolor{dp.cell[1][3]}{state=done}
\recolor{dp.cell[1][4]}{state=current}
\highlight{a.range[0:3]}
\apply{dp.cell[1][4]}{value=8}
\apply{cost_val.cell[0]}{value=8}
\narrate{dp[1][4] = cost(1,4) = 8. Median still at house 3 (cumulative weight 10 $\geq$ 5.5). Cost = 2*2 + 3*1 + 5*0 + 1*1 = 8.}

\step
\recolor{dp.cell[1][4]}{state=done}
\recolor{dp.cell[1][5]}{state=current}
\highlight{a.range[0:4]}
\apply{dp.cell[1][5]}{value=16}
\apply{cost_val.cell[0]}{value=16}
\narrate{dp[1][5] = cost(1,5) = 16. Median at house 3. Cost = 2*2 + 3*1 + 5*0 + 1*1 + 4*2 = 16. Adding distant house 5 increases cost significantly.}

\step
\recolor{dp.cell[1][5]}{state=done}
\recolor{dp.cell[1][6]}{state=current}
\highlight{a.range[0:5]}
\apply{dp.cell[1][6]}{value=33}
\apply{cost_val.cell[0]}{value=33}
\narrate{dp[1][6] = cost(1,6) = 33. Median shifts to house 4. Cost = 2*3 + 3*2 + 5*1 + 1*0 + 4*1 + 6*2 = 33. One school for all houses is expensive!}

\step
\recolor{dp.cell[1][6]}{state=done}
\narrate{Row j=1 complete. With one school, the cost grows quickly as the segment widens. Now we compute row j=2: placing two schools. The recurrence is dp[2][i] = min over m of dp[1][m] + cost(m+1, i).}

\step
\recolor{dp.cell[2][2]}{state=current}
\highlight{a.cell[0]}
\apply{dp.cell[2][2]}{value=0}
\apply{cost_val.cell[0]}{value=0}
\narrate{dp[2][2] = dp[1][1] + cost(2,2) = 0 + 0 = 0. Split at m=1: first school covers house 1, second covers house 2. Both have zero walking.}

\step
\recolor{dp.cell[2][2]}{state=done}
\recolor{dp.cell[2][3]}{state=current}
\highlight{a.range[1:2]}
\apply{cost_val.cell[0]}{value=3}
\narrate{dp[2][3]: try m=1: dp[1][1] + cost(2,3) = 0 + 3 = 3 (median at 3, cost=3*1+5*0=3). Try m=2: dp[1][2] + cost(3,3) = 2 + 0 = 2. Best is m=2.}

\step
\apply{dp.cell[2][3]}{value=2}
\recolor{dp.cell[2][3]}{state=done}
\apply{cost_val.cell[0]}{value=0}
\narrate{dp[2][3] = 2. Optimal split: first school covers houses 1-2 (cost 2), second school at house 3 (cost 0). The D&C optimization exploits that opt(j,i) is monotone.}

\step
\recolor{dp.cell[2][4]}{state=current}
\highlight{a.range[2:3]}
\apply{cost_val.cell[0]}{value=1}
\narrate{dp[2][4]: try m=1 gives 4, m=2 gives dp[1][2]+cost(3,4) = 2+1 = 3 (median at 3, cost=5*0+1*1=1), m=3 gives 7. Best at m=2.}

\step
\apply{dp.cell[2][4]}{value=3}
\recolor{dp.cell[2][4]}{state=done}
\narrate{dp[2][4] = 3. Split at m=2: school 1 covers houses 1-2 (cost 2), school 2 covers 3-4 (cost 1). The monotonicity holds: opt(2,4) = 2 $\geq$ opt(2,3) = 2.}

\step
\recolor{dp.cell[2][5]}{state=current}
\highlight{a.range[3:4]}
\apply{cost_val.cell[0]}{value=1}
\narrate{dp[2][5]: m=3 gives dp[1][3]+cost(4,5) = 7+1 = 8 (school at 5, cost=1*1+4*0=1). m=4 also gives 8+0=8. Best is 8.}

\step
\apply{dp.cell[2][5]}{value=8}
\recolor{dp.cell[2][5]}{state=done}
\narrate{dp[2][5] = 8. With 2 schools for 5 houses, splitting at m=3 or m=4 both yield cost 8. The optimal split point has moved right, confirming monotonicity.}

\step
\recolor{dp.cell[2][6]}{state=current}
\highlight{a.range[4:5]}
\apply{cost_val.cell[0]}{value=4}
\narrate{dp[2][6]: m=4 gives dp[1][4]+cost(5,6) = 8+4 = 12 (school at 6, cost=4*1+6*0=4). Other splits: m=3 gives 13, m=5 gives 16. Best at m=4.}

\step
\apply{dp.cell[2][6]}{value=12}
\recolor{dp.cell[2][6]}{state=done}
\narrate{dp[2][6] = 12. Split at m=4: school 1 at house 3 covers houses 1-4 (cost 8), school 2 at house 6 covers houses 5-6 (cost 4). Total = 12.}

\step
\recolor{dp.cell[2][6]}{state=good}
\recolor{a.cell[2]}{state=good}
\recolor{a.cell[5]}{state=good}
\apply{cost_val.cell[0]}{value=12}
\narrate{Answer: dp[2][6] = 12. Schools at houses 3 and 6. Without D&C optimization this takes O(kn^2). With monotonicity of opt, each row solves in O(n log n), giving O(kn log n) total. For n=3000, that is the difference between TLE and AC.}
\end{animation}
